%% ============================================================
%%  Handout — Aula 6: Gradientes de política II
%%  Curso de Aprendizagem por Reforço — UFU
%%  Prof. Marcelo Keese Albertini
%%  Adaptado e expandido de CS234 (Stanford), Emma Brunskill;
%%  seção de gradientes avançados de Joshua Achiam.
%%  COMPILAR COM XeLaTeX (duas vezes)
%% ============================================================
\documentclass[11pt, a4paper]{article}
%% .sty do projeto na raiz do repositório (../)
\makeatletter
\def\input@path{{./}{../}}
\makeatother


\usepackage{handoutAprendizadoRL}
\usepackage{babel}
\babelprovide[import, main]{brazilian}

\setaula{Aula 6 — Gradientes de política II}
\setprof{Prof. Marcelo Keese Albertini}

%% Macros locais (idênticas às da apresentação)
\newcommand{\Adv}{A}
\newcommand{\Ahat}{\hat{A}}
\newcommand{\ghat}{\hat{g}}
\newcommand{\polnovo}{\pol'}
\newcommand{\dpol}{d^{\pol}}
\newcommand{\dpolnovo}{d^{\pol'}}
\newcommand{\KLd}[2]{D_{\mathrm{KL}}\!\left(#1\,\|\,#2\right)}
\newcommand{\KLbar}{\bar{D}_{\mathrm{KL}}}
\DeclareMathOperator{\clipop}{clip}
\newcommand{\Lsub}[1]{L_{#1}}

\begin{document}

\capa{Aula 6: Gradientes de política II}
     {Linha de base \textbullet\ Ator--crítico \textbullet\ Limites de desempenho relativo \textbullet\ PPO}
     {Prof. Marcelo Keese Albertini}
     {Adaptado de Emma Brunskill, CS234: Reinforcement Learning, Stanford University;
      seção de gradientes avançados a partir dos slides de Joshua Achiam \textbullet\ 2026}

\paragraph{Contexto.} A Aula 5 estabeleceu \emph{para onde} andar: o teorema do
gradiente de política dá uma direção de subida no espaço de parâmetros, e o
REINFORCE a estima sem viés a partir de trajetórias. Esta aula ataca as duas
perguntas que sobraram, e que decidem se o método é utilizável na prática:

\begin{enumerate}
  \item \textbf{Com que precisão conhecemos essa direção?} O estimador do
        REINFORCE é não-viesado, mas sua variância é tão alta que, em lotes de
        tamanho realista, o passo frequentemente aponta para o lado errado.
        A resposta é a \emph{linha de base} e, depois dela, o \emph{crítico}.
  \item \textbf{Qual o maior passo que podemos dar sem destruir a política?}
        Aqui o problema não é estatístico, e sim geométrico: a distância que
        medimos ($\Delta\theta$) não é a distância que importa ($\Delta\pol$).
        A resposta passa por um limite de desempenho relativo entre duas
        políticas e culmina no PPO.
\end{enumerate}

O fio condutor é uma única troca, repetida em três níveis: \textbf{aceitar um
pouco de viés para comprar redução de variância} e, em seguida, \textbf{aceitar
uma aproximação para poder reutilizar dados} --- sempre acompanhada de um
mecanismo explícito que a mantém dentro de sua região de validade.

%% ------------------------------------------------------------
\section{O estimador do REINFORCE e sua variância}

Recordando a Aula 5, com $m$ trajetórias amostradas de $\pol_\theta$:
\[
  \nabla_\theta V(\theta) \;\approx\; \frac{1}{m}\sum_{i=1}^{m}
    \rec(\tau^{(i)}) \sum_{t=0}^{T-1}
    \nabla_\theta \log \pol_\theta\!\left(a_t^{(i)} \mid s_t^{(i)}\right).
\]

\begin{alertabox}
Não-viesado \emph{não} significa utilizável. Um estimador com esperança correta
e variância enorme produz, em cada lote, um passo cuja direção pode ter pouca
relação com o gradiente verdadeiro. O que importa na prática é o erro
quadrático médio a um dado orçamento de amostras.
\end{alertabox}

\subsection{Três fontes de variância, três remédios}

\begin{enumerate}
  \item \textbf{Amostragem da trajetória.} $\tau$ é aleatória: o mesmo $\theta$
        gera retornos diferentes a cada execução. Inevitável --- é o preço de
        não ter modelo.
  \item \textbf{Escala e deslocamento do retorno.} Se todos os retornos valem
        aproximadamente $100$, o estimador empurra \emph{todas} as ações para
        cima; o sinal útil é a diferença entre elas, e essa diferença está
        afogada no nível comum. \emph{Remédio: linha de base.}
  \item \textbf{Atribuição de crédito temporal.} O retorno total $\rec(\tau)$
        multiplica o escore de todas as ações, inclusive as tomadas
        \emph{depois} da recompensa, que não podem tê-la causado.
        \emph{Remédio: estrutura temporal} --- trocar $\rec(\tau)$ por $G_t$
        (feito na Aula 5).
\end{enumerate}

Os itens 2 e 3 são \emph{gratuitos}: reduzem a variância sem introduzir viés.
Um quarto remédio --- trocar $G_t$ por um crítico --- é diferente em natureza:
troca viés por variância deliberadamente, e é o assunto da Seção~\ref{sec:critico}.

%% ------------------------------------------------------------
\section{A linha de base}

\begin{defbox}{Gradiente de política com linha de base}
\[
  \nabla_\theta \E_\tau[\rec] = \E_\tau\!\left[
    \sum_{t=0}^{T-1} \nabla_\theta \log \pol(a_t \mid s_t; \theta)
    \left(\underbrace{\sum_{t'=t}^{T-1} r_{t'}}_{G_t} - b(s_t)\right)
  \right]
\]
onde $b$ é qualquer função \textbf{do estado apenas}. Para toda escolha de $b$
nessas condições, o estimador permanece não-viesado.
\end{defbox}

\begin{ideiabox}
\textbf{Leitura.} Aumente a log-probabilidade de $a_t$ na proporção em que o
retorno obtido foi \emph{melhor do que o esperado naquele estado}. Sem linha de
base, uma ação medíocre tomada num estado excelente é reforçada tanto quanto
uma ação excelente --- porque o que se mede é o retorno absoluto, e não o
mérito da ação.
\end{ideiabox}

\subsection{Demonstração: a linha de base não introduz viés}

\begin{provabox}
Basta mostrar que o termo acrescentado tem esperança nula. Fixe $t$:
\begin{align*}
  \E_\tau&\left[\nabla_\theta \log \pol(a_t \mid s_t;\theta)\, b(s_t)\right] \\
  &= \E_{s_{0:t},\, a_{0:(t-1)}}\Big[\E_{s_{(t+1):T},\, a_{t:(T-1)}}
      \left[\nabla_\theta \log \pol(a_t \mid s_t;\theta)\, b(s_t)\right]\Big]
      && \text{(lei da esperança total)}\\
  &= \E_{s_{0:t},\, a_{0:(t-1)}}\Big[b(s_t)\,
      \E_{a_t\sim\pol(\cdot\mid s_t)}\!\left[\nabla_\theta \log \pol(a_t \mid s_t;\theta)\right]\Big]
      && \text{($b(s_t)$ é constante na esperança interna}\\[-0.35em]
  &&& \text{\phantom{(}e as variáveis após $t$ não aparecem)}\\
  &= \E_{s_{0:t},\, a_{0:(t-1)}}\left[b(s_t)\sum_{a} \pol_\theta(a \mid s_t)
      \frac{\nabla_\theta \pol(a \mid s_t;\theta)}{\pol_\theta(a \mid s_t)}\right]
      && \text{(razão de verossimilhança)}\\
  &= \E_{s_{0:t},\, a_{0:(t-1)}}\left[b(s_t)\, \nabla_\theta \sum_{a} \pol(a \mid s_t;\theta)\right]
      && \text{(soma finita: troca com o gradiente)}\\
  &= \E_{s_{0:t},\, a_{0:(t-1)}}\left[b(s_t)\, \nabla_\theta 1\right] = 0.
      && \text{(normalização de $\pol$)} \qquad\blacksquare
\end{align*}
\end{provabox}

\begin{alertabox}
\textbf{O passo crítico é o terceiro.} Ele exige tirar $b$ de dentro de
$\sum_a$. Se $b$ dependesse da ação --- $b(s,a)$ --- isso seria impossível, o
argumento desmoronaria, e o estimador ficaria viesado. Esta é a pergunta de
prova mais comum sobre o tópico, e a resposta errada mais comum.
\end{alertabox}

\subsection{Por que a variância cai, e qual é a melhor linha de base}

A variância do estimador se decompõe aproximadamente em uma soma sobre passos
de tempo. Concentrando-se em um termo e escrevendo
$\Var[X] = \E[X^2] - (\E[X])^2$, observe que o segundo momento
$(\E[X])^2$ \textbf{não depende de $b$} --- é exatamente o que acabamos de
demonstrar. Minimizar a variância equivale, portanto, a minimizar somente
\[
  \argmin_b \;\E_{s\sim \dpol}\;\E_{a\sim\pol(\cdot\mid s),\, G\mid s,a}
    \left[\left(\nabla_\theta \log \pol(a \mid s;\theta)\right)^2
    \left(G_t(s) - b(s)\right)^2\right].
\]

\begin{teobox}{Linha de base de variância mínima}
Trata-se de um problema de \textbf{mínimos quadrados ponderados}, em que o peso
de cada amostra é o quadrado do escore. Derivando em relação a $b(s)$ e
igualando a zero:
\[
  b^\star(s) \;=\;
  \frac{\E_{a\sim\pol,\,G}\!\left[\left(\nabla_\theta \log \pol(a\mid s;\theta)\right)^2 G_t(s)\right]}
       {\E_{a\sim\pol,\,G}\!\left[\left(\nabla_\theta \log \pol(a\mid s;\theta)\right)^2\right]}
  \;\approx\; \E_{a\sim\pol,\,G}\!\left[G_t(s)\right] \;=\; \Vpi(s).
\]
\end{teobox}

Ou seja: a linha de base ótima é a média dos retornos \emph{ponderada pelo
quadrado da norma do escore}. Ignorando essa ponderação --- o que praticamente
toda implementação faz --- resta a média simples, que é precisamente a
\textbf{função valor de estado} $\Vpi(s)$. É por isso que ``linha de base'' e
``crítico de valor'' acabam sendo a mesma coisa na prática.

\subsection{Exemplo trabalhado: quanto exatamente se ganha?}

\begin{exbox}{resolvido — bandido de três ações}
Considere um problema de um só estado com três ações, política \emph{softmax}
com $\theta = (0,0,0)$ --- portanto $\pol = (\tfrac13,\tfrac13,\tfrac13)$ --- e
recompensas determinísticas $r = (98,\, 101,\, 104)$.

\smallskip
Para a \emph{softmax}, o escore tem forma fechada:
$\nabla_\theta \log \pol(a) = e_a - \pol$, onde $e_a$ é o vetor canônico. Logo
$\norm{\nabla_\theta \log\pol(a)}^2 = 1 - \tfrac13 = \tfrac23$ para toda ação
$a$ --- os pesos do problema de mínimos quadrados são iguais, e portanto
$b^\star$ é a \emph{média simples} dos retornos, $b^\star = 101$.

\smallskip
O estimador de uma única amostra é $\ghat = (e_a - \pol)(r(a) - b)$. Calculando
a esperança e a variância total (traço da matriz de covariância):

\medskip
\begin{center}
\begin{tabular}{@{}lccc@{}}
\toprule
 & $\E[\ghat]$ & $\E\norm{\ghat}^2$ & $\Var$ total \\
\midrule
sem linha de base ($b=0$)     & $(-1,\,0,\,+1)$ & $6804{,}67$ & $\mathbf{6802{,}67}$ \\
com linha de base ($b=101$)   & $(-1,\,0,\,+1)$ & $4{,}00$    & $\mathbf{2{,}00}$ \\
\bottomrule
\end{tabular}
\end{center}
\vspace{-0.3em}
O gradiente esperado é \textbf{idêntico} nos dois casos, como a demonstração
garante. A variância cai por um fator de $\approx 3\,400$. Traduzindo para
orçamento de amostras: como o erro-padrão da média decai com $\sqrt{m}$, seriam
necessárias cerca de $3\,400$ vezes mais trajetórias para o estimador sem linha
de base atingir a mesma precisão.
\end{exbox}

\begin{ideiabox}
O exemplo é extremo de propósito, mas o mecanismo é geral: o ganho cresce com a
razão entre o \emph{nível} dos retornos e a sua \emph{dispersão}. Ambientes com
recompensa de sobrevivência (``$+1$ por passo vivo''), muito comuns, estão
exatamente nesse regime.
\end{ideiabox}

%% ------------------------------------------------------------
\section{Do retorno ao crítico: métodos ator--crítico}
\label{sec:critico}

O retorno $G_t^i$ é uma estimativa não-viesada de $\Vpi(s_t)$ obtida de uma
\emph{única} execução --- soma de muitas variáveis aleatórias, com variância
proporcional ao horizonte. A alternativa é reduzir essa variância
\textbf{introduzindo viés}, via \emph{bootstrapping} e aproximação de função.
É exatamente o dilema TD versus Monte Carlo da Aula 3, agora aplicado ao alvo
do gradiente de política.

\begin{defbox}{Métodos ator--crítico}
Mantêm representações explícitas \emph{da política} (o \textbf{ator},
$\pol_\theta$) e \emph{da função valor} (o \textbf{crítico}, $V_w$), e
atualizam as duas em paralelo. A estimativa de $V$ ou $Q$ que entra no
gradiente de política é produzida pelo crítico, e não pelo retorno observado.
\end{defbox}

\subsection{A função vantagem}

Partindo da forma com linha de base e substituindo o retorno amostral por um
crítico $Q(s_t,a_t;w)$, e depois tomando a linha de base como uma estimativa de
$V$, a expressão colapsa em:
\[
  \nabla_\theta \E_\tau[\rec] \;\approx\; \E_\tau\!\left[\sum_{t=0}^{T-1}
    \nabla_\theta \log \pol(a_t\mid s_t;\theta)\,
    \mdestaque{rlLaranja}{\Ahat^{\pol}(s_t,a_t)}\right],
  \qquad
  \Adv^{\pol}(s,a) = \Qpi(s,a) - \Vpi(s).
\]

\begin{teobox}{Propriedade fundamental da vantagem}
$\E_{a\sim\pol(\cdot\mid s)}\!\left[\Adv^{\pol}(s,a)\right] = 0$ para todo $s$.
\end{teobox}

\begin{provabox}
Imediato da definição de $\Vpi$ como esperança de $\Qpi$ sob a política:
$\E_{a\sim\pol}[\Adv^{\pol}(s,a)] = \E_{a\sim\pol}[\Qpi(s,a)] - \Vpi(s)
= \Vpi(s) - \Vpi(s) = 0$. \hfill$\blacksquare$
\end{provabox}

\begin{ideiabox}
A vantagem é, por construção, \textbf{centrada}: metade das ações da política
atual tem vantagem positiva e metade negativa, em média ponderada. Essa é a
razão estatística de sua eficácia --- é a versão ``com média zero'' do sinal de
aprendizado, e não apenas uma quantidade conveniente.
\end{ideiabox}

\subsection{O catálogo de alvos}

Todos os pesos abaixo podem substituir $G_t$ no gradiente. Descer na tabela é
caminhar do Monte Carlo puro para o TD puro.

\medskip
\begin{center}
\begin{tabular}{@{}llll@{}}
\toprule
\textbf{Peso do escore} & \textbf{Nome} & \textbf{Viés} & \textbf{Variância} \\
\midrule
$\rec(\tau)$ & retorno total da trajetória & nenhum & altíssima \\
$G_t = \sum_{t'\ge t} r_{t'}$ & retorno futuro & nenhum & muito alta \\
$G_t - b(s_t)$ & retorno com linha de base & nenhum & alta \\
$G_t - V_w(s_t)$ & vantagem de Monte Carlo & só de $V_w$ & média \\
$Q_w(s_t,a_t) - V_w(s_t)$ & vantagem estimada & de $Q_w$ e $V_w$ & baixa \\
$\delta_t = r_t + \gamma V_w(s_{t+1}) - V_w(s_t)$ & erro de TD & maior & muito baixa \\
$\Ahat_t^{\mathrm{GAE}(\gamma,\lambda)}$ & vantagem generalizada & ajustável & ajustável \\
\bottomrule
\end{tabular}
\end{center}
\vspace{-0.3em}
A última linha interpola continuamente entre as demais por meio de um único
parâmetro $\lambda \in [0,1]$, e é o assunto da Aula 7. Note que
$\E[\delta_t \mid s_t,a_t] = \Adv^{\pol}(s_t,a_t)$ \emph{apenas} quando
$V_w = \Vpi$; fora disso, o erro de TD é um estimador viesado da vantagem, e
esse viés é o preço explícito da redução de variância.

%% ------------------------------------------------------------
\section{Os problemas estruturais do gradiente de política}

O objetivo é
$\max_\theta J(\pol_\theta) \defeq \E_{\tau\sim\pol_\theta}\big[\sum_{t\ge 0}\gamma^t r_t\big]$,
atacado por subida de gradiente estocástica com
$g = \E_{\tau\sim\pol_\theta}\big[\sum_{t\ge 0}\gamma^t
\nabla_\theta \log \pol_\theta(a_t\mid s_t)\Adv^{\pol_\theta}(s_t,a_t)\big]$.
Três limitações persistem mesmo com variância controlada.

\subsection{Eficiência amostral}

O gradiente de política é uma esperança \textbf{sob a política atual}: é
\emph{on-policy} por construção. Assim que $\theta$ muda, os dados coletados
descrevem uma política que não existe mais, e o lote inteiro precisa ser
descartado após um único passo de gradiente.

\begin{ideiabox}
\textbf{Oportunidade.} Reaproveitar dados antigos para dar \emph{vários} passos
antes de coletar de novo. \textbf{Desafio.} Quantos passos podemos dar antes que
os dados deixem de descrever a política que estamos otimizando? Esta é,
literalmente, a pergunta que o PPO responde.
\end{ideiabox}

\subsection{Espaço de parâmetros \emph{versus} espaço de políticas}

\begin{defbox}{Espaço de políticas}
No caso tabular, o conjunto de matrizes estocásticas por linha
$\Pi = \left\{\pol \in \R^{\abs{\gS}\times\abs{\gA}} :
\sum_a \pol_{sa} = 1,\; \pol_{sa}\ge 0\right\}$.
O gradiente dá passos em $\theta$; o desempenho depende do deslocamento em
$\Pi$. Não há razão para que a mesma régua sirva nos dois lugares.
\end{defbox}

\begin{exbox}{resolvido — a política sigmoide}
Seja $\pol_\theta(a{=}1) = \sigma(\theta)$ e $\pol_\theta(a{=}2) = 1-\sigma(\theta)$,
com um único parâmetro escalar. Compare dois passos de \emph{mesmo tamanho} em
$\theta$:

\medskip
\begin{center}
\begin{tabular}{@{}lcccc@{}}
\toprule
Passo & $\pol$ inicial & $\pol$ final & $\Delta\pol$ & $\KLd{\pol_{\text{final}}}{\pol_{\text{inicial}}}$ \\
\midrule
$\theta: -1 \to +1$ \; ($\Delta\theta = 2$) & $0{,}2689$ & $0{,}7311$ & $\mathbf{0{,}4621}$ & $\mathbf{0{,}4621}$ \\
$\theta: \;\;4 \to \;\,6$ \; ($\Delta\theta = 2$) & $0{,}9820$ & $0{,}9975$ & $\mathbf{0{,}0155}$ & $\mathbf{0{,}0107}$ \\
\bottomrule
\end{tabular}
\end{center}
\vspace{-0.3em}
O mesmo $\Delta\theta = 2$ produz mudanças de comportamento que diferem por um
fator de $30$ em probabilidade e de $43$ em divergência KL. Um passo calibrado
para a primeira região é catastrófico ou inútil na segunda, conforme a direção.
\end{exbox}

\begin{alertabox}
\textbf{Consequência prática: o colapso de desempenho.} Um passo grande demais
pode derrubar $J$ para um patamar do qual não se sai --- porque a política ruim
coleta dados ruins, com os quais não se aprende a escapar. Passo pequeno demais
significa progresso inaceitavelmente lento. E o tamanho ``certo'' \emph{muda com
$\theta$}: nenhuma constante serve para a corrida inteira.
\end{alertabox}

\subsection{Uma observação: KL como métrica local}

A ligação entre as duas colunas do exemplo anterior não é acidental. Para
$\polnovo$ próxima de $\pol$, a expansão de Taylor de segunda ordem da
divergência KL é
\[
  \KLbar(\theta + \Delta\theta \,\|\, \theta)
  \;\approx\; \tfrac{1}{2}\,\Delta\theta^\top F(\theta)\, \Delta\theta,
  \qquad
  F(\theta) = \E_{s\sim\dpol,\, a\sim\pol}\!\left[
    \nabla_\theta \log \pol_\theta(a\mid s)\,
    \nabla_\theta \log \pol_\theta(a\mid s)^\top\right],
\]
onde $F$ é a \textbf{matriz de informação de Fisher}. Ou seja: a KL é
localmente uma norma quadrática cuja métrica depende de $\theta$. Restringir a
KL é medir o passo com a régua correta em cada ponto --- e usar $F^{-1}g$ em
vez de $g$ como direção é precisamente o \emph{gradiente natural}, sobre o qual
o TRPO se constrói. O PPO evita calcular $F$, e paga por isso em rigor.

%% ------------------------------------------------------------
\section{Limites de desempenho relativo entre duas políticas}

Queremos uma regra de atualização que (i) use as trajetórias da política mais
recente com a maior eficiência possível e (ii) dê passos que respeitem
distância no espaço de políticas. Para isso precisamos relacionar o desempenho
de \emph{duas} políticas.

\begin{teobox}{Lema da diferença de desempenho}
Para quaisquer políticas $\pol$ e $\polnovo$,
\[
  J(\polnovo) - J(\pol)
  \;=\; \E_{\tau\sim\polnovo}\!\left[\sum_{t=0}^{\infty}\gamma^t \Adv^{\pol}(s_t,a_t)\right]
  \;=\; \frac{1}{1-\gamma}\,
        \E_{\substack{s\sim\dpolnovo\\ a\sim\polnovo}}\!\left[\Adv^{\pol}(s,a)\right],
\]
onde $\dpol(s) = (1-\gamma)\sum_{t\ge 0}\gamma^t\,\Prob(s_t = s \mid \pol)$ é a
\textbf{distribuição descontada de estados futuros}.
\end{teobox}

\begin{provabox}
Escreva a vantagem como um resíduo de Bellman:
$\Adv^{\pol}(s,a) = \E_{s'}\!\left[\rec(s,a) + \gamma \Vpi(s') - \Vpi(s)\right]$.
Então
\begin{align*}
  \E_{\tau\sim\polnovo}\!\left[\sum_{t\ge 0}\gamma^t \Adv^{\pol}(s_t,a_t)\right]
  &= \E_{\tau\sim\polnovo}\!\left[\sum_{t\ge 0}\gamma^t r_t\right]
   + \E_{\tau\sim\polnovo}\!\left[\sum_{t\ge 0}
     \left(\gamma^{t+1}\Vpi(s_{t+1}) - \gamma^{t}\Vpi(s_t)\right)\right] \\
  &= J(\polnovo) \;+\; \E_{s_0}\!\left[-\Vpi(s_0)\right]
   \;=\; J(\polnovo) - J(\pol),
\end{align*}
pois a segunda soma é \textbf{telescópica} e sobra apenas o termo $t=0$ com
sinal negativo. A segunda igualdade do enunciado é a definição de $\dpolnovo$:
trocar $\sum_t \gamma^t \E_{s_t}[\,\cdot\,]$ por
$\frac{1}{1-\gamma}\E_{s\sim\dpolnovo}[\,\cdot\,]$. \hfill$\blacksquare$
\end{provabox}

\begin{ideiabox}
\textbf{Leitura.} A melhoria de $\polnovo$ sobre $\pol$ é a vantagem média,
medida com a \emph{régua antiga} $\Adv^{\pol}$, ao longo das \emph{trajetórias
novas}. Toda a teoria que segue é uma tentativa de tornar esse lado direito
computável com dados que já temos.
\end{ideiabox}

\subsection{Amostragem por importância e o obstáculo restante}

Reescrevendo a esperança sobre ações com um peso de amostragem por importância:
\[
  J(\polnovo) - J(\pol)
  = \frac{1}{1-\gamma}\E_{\substack{s\sim\dpolnovo\\ a\sim\polnovo}}\!\left[\Adv^{\pol}(s,a)\right]
  = \frac{1}{1-\gamma}\E_{\substack{s\sim\dpolnovo\\ a\sim\pol}}
    \!\left[\frac{\polnovo(a\mid s)}{\pol(a\mid s)}\,\Adv^{\pol}(s,a)\right].
\]

As \emph{ações} já vêm da política antiga. O obstáculo que resta é
$s \sim \dpolnovo$: os \textbf{estados} ainda são visitados pela política nova,
que é justamente a que ainda não temos.

\subsection{A aproximação e o limite}

E se simplesmente disséssemos $\dpolnovo \approx \dpol$?
\[
  J(\polnovo) - J(\pol) \;\approx\;
  \frac{1}{1-\gamma}\E_{\substack{s\sim\dpol\\ a\sim\pol}}\!\left[
  \frac{\polnovo(a\mid s)}{\pol(a\mid s)}\,\Adv^{\pol}(s,a)\right]
  \;\defeq\; \Lsub{\pol}(\polnovo).
\]

\begin{teobox}{Limite de desempenho relativo \normalfont (Achiam, Held, Tamar e Abbeel, 2017)}
\[
  \abs{\,J(\polnovo) - \left(J(\pol) + \Lsub{\pol}(\polnovo)\right)}
  \;\le\; C\,\sqrt{\E_{s\sim\dpol}\!\left[\KLd{\polnovo}{\pol}[s]\right]}.
\]
\end{teobox}

O erro cometido ao trocar $\dpolnovo$ por $\dpol$ é controlado pela divergência
KL média entre as políticas. Se elas são próximas em KL, a aproximação é boa e
$\Lsub{\pol}(\polnovo)$ torna-se um \textbf{substituto} legítimo do objetivo.

\begin{defbox}{Divergência de Kullback--Leibler}
$\KLd{P}{Q} = \sum_x P(x)\log\frac{P(x)}{Q(x)}$, e entre políticas num estado
$s$: $\KLd{\polnovo}{\pol}[s] = \sum_{a\in\gA}\polnovo(a\mid s)\log
\frac{\polnovo(a\mid s)}{\pol(a\mid s)}$.
Propriedades: $\KLd{P}{P} = 0$; \; $\KLd{P}{Q}\ge 0$; \;
$\KLd{P}{Q}\ne\KLd{Q}{P}$ --- \textbf{não é simétrica}, logo não é uma métrica.
\end{defbox}

\subsection{O que se ganha com a aproximação}

\[
  \Lsub{\pol}(\polnovo)
  = \E_{\tau\sim\pol}\!\left[\sum_{t=0}^{\infty}\gamma^t
    \frac{\polnovo(a_t\mid s_t)}{\pol(a_t\mid s_t)}\,\Adv^{\pol}(s_t,a_t)\right].
\]

Tudo do lado direito é estimável com trajetórias amostradas da política
\textbf{antiga}. Podemos otimizar $\polnovo$ sem coletar dados novos --- isto
é, dar \emph{vários passos de gradiente por lote}.

\begin{ideiabox}
\textbf{Por que os pesos não explodem.} Na amostragem por importância aplicada
a trajetórias inteiras, o peso é um \emph{produto} de $t$ razões,
$\prod_{t'=0}^{t}\frac{\polnovo(a_{t'}\mid s_{t'})}{\pol(a_{t'}\mid s_{t'})}$;
mesmo diferenças minúsculas se acumulam multiplicativamente e o estimador
degenera. Em $\Lsub{\pol}$ o peso depende \textbf{só do passo atual}. Não há
produto, e portanto não há explosão nem desaparecimento.
\end{ideiabox}

\begin{alertabox}
A licença é condicional. $\Lsub{\pol}(\polnovo)$ só aproxima o objetivo
verdadeiro \emph{enquanto $\polnovo$ permanecer próxima de $\pol$ em KL}.
Maximizar $\Lsub{\pol}$ sem restrição alguma é otimizar um modelo fora de sua
região de validade --- e o máximo irrestrito costuma ser exatamente o lugar
onde o modelo é menos confiável.
\end{alertabox}

%% ------------------------------------------------------------
\section{Otimização proximal de política (PPO)}

O PPO é uma família de métodos que impõe a restrição de KL \textbf{de forma
aproximada}, sem calcular gradientes naturais. Duas variantes.

\subsection{Variante 1 --- penalidade KL adaptativa}

A atualização resolve um problema \emph{sem restrição}, em que a KL entra como
penalidade:
\[
  \theta_{k+1} = \argmax_\theta\; \Lsub{\theta_k}(\theta)
    - \beta_k\, \KLbar(\theta \| \theta_k),
  \qquad
  \KLbar(\theta\|\theta_k) = \E_{s\sim d^{\pol_k}}\!\left[
    \KLd{\pol_{\theta_k}(\cdot\mid s)}{\pol_\theta(\cdot\mid s)}\right].
\]

\begin{algbox}{PPO com penalidade KL adaptativa}
\textbf{Entrada:} $\theta_0$; penalidade inicial $\beta_0$; KL alvo $\delta$.
\begin{enumerate}[leftmargin=1.6em, itemsep=0.1em]
  \item \textbf{para} $k = 0,1,2,\dots$:
  \begin{enumerate}[leftmargin=1.4em, itemsep=0.05em, label=\alph*)]
    \item Colete trajetórias parciais $\mathcal{D}_k$ executando $\pol_k = \pol(\theta_k)$.
    \item Estime as vantagens $\Ahat_t^{\pol_k}$ por qualquer método.
    \item Calcule $\theta_{k+1} = \argmax_\theta \Lsub{\theta_k}(\theta)
          - \beta_k \KLbar(\theta\|\theta_k)$ com $K$ passos de SGD por
          minilotes (Adam).
    \item \textbf{se} $\KLbar(\theta_{k+1}\|\theta_k) \ge 1{,}5\,\delta$
          \textbf{então} $\beta_{k+1} \leftarrow 2\beta_k$; \\
          \textbf{senão se} $\KLbar(\theta_{k+1}\|\theta_k) \le \delta/1{,}5$
          \textbf{então} $\beta_{k+1} \leftarrow \beta_k/2$.
  \end{enumerate}
\end{enumerate}
\end{algbox}

O valor inicial de $\beta$ importa pouco --- ele se adapta rapidamente. Algumas
iterações violam a restrição; a maioria não. O parâmetro decisivo é $K$: é aqui
que finalmente damos \emph{vários} passos de gradiente por lote de dados.

\subsection{Variante 2 --- objetivo recortado}

\begin{defbox}{Razão de probabilidades e objetivo $L^{\mathrm{CLIP}}$}
Seja $r_t(\theta) = \dfrac{\pol_\theta(a_t\mid s_t)}{\pol_{\theta_k}(a_t\mid s_t)}$. Então
\[
  L^{\mathrm{CLIP}}_{\theta_k}(\theta) = \E_{\tau\sim\pol_k}\!\left[
    \sum_{t=0}^{T}
    \min\!\Big(r_t(\theta)\,\Ahat_t,\;
    \clipop\big(r_t(\theta),\, 1-\epsilon,\, 1+\epsilon\big)\,\Ahat_t\Big)\right],
\]
com $\epsilon$ hiperparâmetro (tipicamente $0{,}2$), e
$\theta_{k+1} = \argmax_\theta L^{\mathrm{CLIP}}_{\theta_k}(\theta)$.
\end{defbox}

No início de cada época de otimização temos $\theta = \theta_k$, logo
$r_t = 1$ e $L^{\mathrm{CLIP}} = \Lsub{\theta_k}$: o objetivo recortado coincide
com o substituto exatamente onde o substituto é confiável.

\subsection{Análise caso a caso do recorte}

Desenvolvendo o $\min$ nas seis combinações possíveis:

\medskip
\begin{center}
\begin{tabular}{@{}lccc@{}}
\toprule
 & $r_t < 1-\epsilon$ & $1-\epsilon \le r_t \le 1+\epsilon$ & $r_t > 1+\epsilon$ \\
\midrule
$\Ahat_t > 0$ (ação boa)  & $r_t\Ahat_t$ \; \emph{ativo} & $r_t\Ahat_t$ \; \emph{ativo}
  & $(1+\epsilon)\Ahat_t$ \; \textbf{plano} \\
$\Ahat_t < 0$ (ação ruim) & $(1-\epsilon)\Ahat_t$ \; \textbf{plano} & $r_t\Ahat_t$ \; \emph{ativo}
  & $r_t\Ahat_t$ \; \emph{ativo} \\
\bottomrule
\end{tabular}
\end{center}
\vspace{-0.3em}
``Plano'' significa gradiente nulo em relação a $\theta$: o otimizador deixa de
empurrar naquela direção.

\begin{ideiabox}
\textbf{A assimetria é deliberada.} O $\min$ escolhe sempre o ramo mais
\emph{pessimista}. Para uma ação boa, o recorte remove o incentivo de aumentar
sua probabilidade além de $1+\epsilon$. Para uma ação ruim que já foi empurrada
para $r_t \gg 1$ por um passo anterior, o gradiente \textbf{continua ativo} ---
o recorte nunca remove o incentivo de \emph{voltar}. É exatamente o
comportamento que se deseja de um mecanismo de segurança.
\end{ideiabox}

\begin{alertabox}
O recorte limita a \textbf{razão}, não a divergência KL, e o limite é por
amostra, não em média. Não há garantia formal de melhoria monotônica: o PPO
recortado é uma heurística barata \emph{inspirada} na teoria da seção anterior,
não uma consequência dela. Na prática, deve-se monitorar a KL empírica.
\end{alertabox}

\subsection{O algoritmo completo, na forma usada em implementações}

\begin{algbox}{PPO com objetivo recortado}
\textbf{Entrada:} $\theta_0$ (ator), $w_0$ (crítico), $\epsilon$, épocas $K$,
tamanho de minilote.
\begin{enumerate}[leftmargin=1.6em, itemsep=0.1em]
  \item \textbf{para} $k = 0,1,2,\dots$:
  \begin{enumerate}[leftmargin=1.4em, itemsep=0.05em, label=\alph*)]
    \item Execute $\pol_{\theta_k}$ por $T$ passos em $N$ ambientes paralelos.
    \item Estime as vantagens $\Ahat_t$ (na prática, por GAE --- Aula 7) e
          \textbf{congele-as}.
    \item Guarde $\log \pol_{\theta_k}(a_t\mid s_t)$: é o denominador de $r_t(\theta)$.
    \item \textbf{para} época $=1,\dots,K$, e para cada minilote: um passo de
          Adam sobre
          $\;-L^{\mathrm{CLIP}} + c_1 L^{V} - c_2\,\mathcal{H}[\pol_\theta]$.
  \end{enumerate}
\end{enumerate}
\end{algbox}

São três termos na perda: o objetivo recortado do \textbf{ator}; o erro
quadrático do \textbf{crítico}, $L^{V} = (V_w(s_t) - \hat{G}_t)^2$; e um bônus
de \textbf{entropia} $\mathcal{H}$, que impede a política de se tornar
determinística cedo demais e matar a exploração.

\subsection{Valores usuais e armadilhas}

\begin{itemize}
  \item $\epsilon \approx 0{,}2$; \; $K$ entre $3$ e $10$ épocas por lote;
        $c_1 \approx 0{,}5$; \; $c_2 \approx 0{,}01$.
  \item Vantagens normalizadas por minilote; corte da norma do gradiente em
        $0{,}5$.
  \item \textbf{Monitorar} a KL empírica entre $\pol_\theta$ e $\pol_{\theta_k}$
        e a \emph{fração de amostras recortadas}. KL disparando $\Rightarrow$
        reduzir $K$ ou a taxa de aprendizado.
\end{itemize}

\begin{alertabox}
\textbf{Erros de implementação frequentes.}
\begin{itemize}[itemsep=0.1em]
  \item \emph{Recalcular} $\Ahat_t$ dentro do laço de épocas. As vantagens são
        de $\pol_{\theta_k}$ e devem permanecer congeladas durante as $K$ épocas.
  \item Não destacar o gradiente do denominador de $r_t(\theta)$: ele é uma
        constante, não uma função de $\theta$.
  \item $K$ grande demais. A política se afasta, e o recorte sozinho não segura
        a KL --- é a falha mais comum e a mais difícil de diagnosticar, porque
        a curva de aprendizado piora lentamente em vez de quebrar.
\end{itemize}
Boa parte do desempenho publicado do PPO vem dessas escolhas de implementação,
e não do objetivo em si.
\end{alertabox}

%% ------------------------------------------------------------
\section{Formulário}

\begin{tcolorbox}[enhanced, colback=rlFundo, colframe=rlAzul!25, boxrule=0.8pt,
  arc=1.6mm, left=3mm, right=3mm, top=2.4mm, bottom=2.4mm, breakable]
\textbf{1. Gradiente de política com linha de base}
\[
  \nabla_\theta J = \E_\tau\!\left[\sum_{t}
    \nabla_\theta \log \pol_\theta(a_t\mid s_t)\left(G_t - b(s_t)\right)\right],
  \qquad
  \E\!\left[\nabla_\theta \log\pol_\theta(a\mid s)\, b(s)\right] = 0.
\]

\textbf{2. Linha de base ótima e vantagem}
\[
  b^\star(s) \approx \Vpi(s),
  \qquad
  \Adv^{\pol}(s,a) = \Qpi(s,a) - \Vpi(s),
  \qquad
  \E_{a\sim\pol}\!\left[\Adv^{\pol}(s,a)\right] = 0.
\]

\textbf{3. Distribuição descontada de estados e lema da diferença de desempenho}
\[
  \dpol(s) = (1-\gamma)\sum_{t\ge0}\gamma^t\Prob(s_t = s\mid\pol),
  \qquad
  J(\polnovo) - J(\pol) = \frac{1}{1-\gamma}
  \E_{\substack{s\sim\dpolnovo\\a\sim\polnovo}}\!\left[\Adv^{\pol}(s,a)\right].
\]

\textbf{4. Substituto e limite de desempenho relativo}
\[
  \Lsub{\pol}(\polnovo) = \frac{1}{1-\gamma}
  \E_{\substack{s\sim\dpol\\a\sim\pol}}\!\left[
  \frac{\polnovo(a\mid s)}{\pol(a\mid s)}\Adv^{\pol}(s,a)\right]
  = \E_{\tau\sim\pol}\!\left[\sum_{t\ge0}\gamma^t
  \frac{\polnovo(a_t\mid s_t)}{\pol(a_t\mid s_t)}\Adv^{\pol}(s_t,a_t)\right],
\]
\[
  \abs{J(\polnovo) - J(\pol) - \Lsub{\pol}(\polnovo)}
  \;\le\; C\sqrt{\E_{s\sim\dpol}\!\left[\KLd{\polnovo}{\pol}[s]\right]}.
\]

\textbf{5. PPO}
\[
  r_t(\theta) = \frac{\pol_\theta(a_t\mid s_t)}{\pol_{\theta_k}(a_t\mid s_t)},
  \qquad
  L^{\mathrm{CLIP}} = \E\!\left[\sum_t \min\!\Big(r_t\Ahat_t,\;
  \clipop(r_t, 1-\epsilon, 1+\epsilon)\,\Ahat_t\Big)\right].
\]

\textbf{6. KL local e informação de Fisher}
\[
  \KLbar(\theta+\Delta\theta\,\|\,\theta) \approx
  \tfrac12 \Delta\theta^\top F(\theta)\Delta\theta,
  \qquad
  F(\theta) = \E\!\left[\nabla_\theta\log\pol_\theta\,\nabla_\theta\log\pol_\theta^\top\right].
\]
\end{tcolorbox}

%% ------------------------------------------------------------
\section{Exercícios}

\begin{exbox}{1 — a linha de base precisa mesmo ser só do estado?}
Considere um bandido de duas ações com $\pol = (p,\, 1-p)$ e recompensas
$r_1, r_2$. Tome a ``linha de base'' dependente da ação
$b(a) = r(a)$, que zera todo o peso do escore.
\begin{enumerate}[label=(\alph*)]
  \item Calcule o gradiente esperado do estimador resultante.
  \item Compare-o com o gradiente verdadeiro. O estimador é viesado?
  \item Aponte o passo exato da demonstração da Seção~2 que falha.
\end{enumerate}
\end{exbox}

\begin{exbox}{2 — deslocamento constante das recompensas}
Suponha que todas as recompensas do problema sejam deslocadas por uma constante
$c$: $r'(s,a) = r(s,a) + c$, em episódios de comprimento fixo $T$.
\begin{enumerate}[label=(\alph*)]
  \item Mostre que o gradiente \emph{esperado} do REINFORCE não muda.
  \item Mostre que a variância do estimador de uma amostra cresce com $c^2$.
  \item Explique por que introduzir $b(s) = \E[G_t]$ elimina completamente essa
        dependência em $c$.
\end{enumerate}
\end{exbox}

\begin{exbox}{3 — o erro de TD como estimador da vantagem}
Seja $\delta_t = r_t + \gamma V_w(s_{t+1}) - V_w(s_t)$.
\begin{enumerate}[label=(\alph*)]
  \item Mostre que, se $V_w = \Vpi$, então
        $\E\!\left[\delta_t \mid s_t = s,\, a_t = a\right] = \Adv^{\pol}(s,a)$.
  \item Escreva o viés de $\delta_t$ como estimador da vantagem em função do
        erro $\varepsilon(s) = V_w(s) - \Vpi(s)$.
  \item Em que condição sobre $\varepsilon$ esse viés se anula mesmo com
        $V_w \ne \Vpi$?
\end{enumerate}
\end{exbox}

\begin{exbox}{4 — o substituto estende o gradiente de política}
Seja $\pol_{\theta'}$ a política nova e $\pol_\theta$ a antiga.
\begin{enumerate}[label=(\alph*)]
  \item Mostre que $\Lsub{\pol_\theta}(\pol_\theta) = 0$.
  \item Mostre que $\nabla_{\theta'}\Lsub{\pol_\theta}(\pol_{\theta'})
        \big|_{\theta'=\theta}$ é exatamente o gradiente de política da Aula 5.
  \item Explique, a partir de (b), por que otimizar $\Lsub{\pol}$ por um único
        passo pequeno é equivalente a um passo de gradiente de política comum
        --- e o que se ganha ao dar vários passos.
\end{enumerate}
\end{exbox}

\begin{exbox}{5 — derivando a tabela do recorte}
Sem consultar a tabela da Seção~6, deduza os seis casos de
$\min\!\left(r\Ahat,\; \clipop(r,1-\epsilon,1+\epsilon)\Ahat\right)$ conforme o
sinal de $\Ahat$ e a posição de $r$. Em seguida esboce as duas curvas
$L^{\mathrm{CLIP}}(r)$ e identifique, em cada uma, onde o gradiente é nulo.
\end{exbox}

\begin{exbox}{6 — calibrando o passo pela KL}
Para a política sigmoide $\pol_\theta(a{=}1) = \sigma(\theta)$:
\begin{enumerate}[label=(\alph*)]
  \item Calcule numericamente o maior $\Delta\theta > 0$ tal que
        $\KLd{\pol_{\theta+\Delta\theta}}{\pol_\theta} \le 0{,}01$, para
        $\theta = 0$ e para $\theta = 4$.
  \item Compare os dois valores. Qual seria o efeito de usar um passo constante
        calibrado em $\theta = 0$ quando a política já chegou a $\theta = 4$?
  \item Verifique a aproximação de Fisher da Seção~4 comparando o resultado
        exato com $\Delta\theta \approx \sqrt{2\cdot 0{,}01 / F(\theta)}$, onde
        $F(\theta) = \sigma(\theta)(1-\sigma(\theta))$ para esta família.
\end{enumerate}
\end{exbox}

\begin{exbox}{7 — implementação}
Na implementação de PPO da lista de exercícios, produza gráficos, ao longo do
treinamento, de: (i) a KL empírica entre $\pol_\theta$ e $\pol_{\theta_k}$ ao
fim de cada iteração; (ii) a fração de amostras em que o recorte foi ativado;
(iii) o retorno médio. Rode com $K \in \{1,\, 4,\, 20\}$ épocas e relacione o
que se vê em (i) e (ii) com o que acontece em (iii).
\end{exbox}

%% ------------------------------------------------------------
\needspace{0.42\textheight}
\section{Para discutir em aula}

\begin{discbox}
\begin{enumerate}[itemsep=0.35em]
  \item Uma linha de base que dependa da \emph{ação} preservaria o não-viés do
        estimador? Onde exatamente a demonstração quebraria?
  \item Trocar $G_t - b(s_t)$ por $\delta_t$ preserva o não-viés? Reduz a
        variância? As duas respostas são independentes uma da outra?
  \item Se $\Lsub{\pol}(\polnovo) > 0$, podemos concluir que
        $J(\polnovo) > J(\pol)$? E se soubermos, além disso, que a KL média é
        pequena?
  \item O limite de desempenho relativo deixa de \emph{valer} quando as
        políticas estão distantes, ou apenas de ser \emph{informativo}?
  \item Nos dois gráficos de $L^{\mathrm{CLIP}}$, qual corresponde a
        $\Adv > 0$ e qual a $\Adv < 0$? A resposta depende de $\epsilon$?
  \item Por que o recorte é assimétrico --- por que não achatar a curva dos
        dois lados? O que se perderia?
  \item O PPO garante melhoria monotônica? Se não, o que herda da Seção~5?
\end{enumerate}
\end{discbox}

%% ------------------------------------------------------------
\section*{Referências}
\begin{itemize}[itemsep=0.15em]
  \item S.~Kakade e J.~Langford. \emph{Approximately Optimal Approximate
        Reinforcement Learning}. ICML, 2002.
  \item J.~Schulman, S.~Levine, P.~Moritz, M.~Jordan e P.~Abbeel.
        \emph{Trust Region Policy Optimization}. ICML, 2015.
  \item J.~Achiam, D.~Held, A.~Tamar e P.~Abbeel. \emph{Constrained Policy
        Optimization}. ICML, 2017.
  \item J.~Schulman, F.~Wolski, P.~Dhariwal, A.~Radford e O.~Klimov.
        \emph{Proximal Policy Optimization Algorithms}. arXiv:1707.06347, 2017.
  \item V.~Mnih et al. \emph{Asynchronous Methods for Deep Reinforcement
        Learning} (A3C). ICML, 2016.
\end{itemize}

\vspace{0.6em}
{\footnotesize\color{rlCinza}Material adaptado e expandido de \emph{CS234:
Reinforcement Learning}, Emma Brunskill, Stanford University; a seção de
gradientes de política avançados segue os slides de Joshua Achiam.}

\end{document}
